\section{东方封建网络：不是地图上的颜色，而是一套交换关系}

周人在东方设立封国，常被一句“分封制”概括。但在当时，封国不是中央任命的普通地方政府：它们有宗族、军队、祭祀和本地关系，也对周王承担朝聘、贡赋、军事与礼仪义务。网络的关键不在于周王是否能每天发号施令，而在于冲突发生时，各方是否仍愿意承认周王室有分配名分和裁决关系的资格。

\begin{event}{西周早期}{营建成周与经营东方}{年代有争议}{洛邑及东方地区}\eventanchor{WZ-EAST-CHENGZHOU}{西周·营建成周}
\term{这一阶段发生了什么}：传统材料将成周洛邑的营建、殷遗民的迁置和东方封国的重新安排联系在一起。洛邑既靠近东方，也能成为周王室向东展开礼仪与军事影响的支点。

\term{事情为什么重要}：牧野之后，周若只守住关中，便无法接住商的政治遗产。成周与封国网络降低了王室直接驻军的成本，却把治理变成持续的关系维护：册命、婚姻、赏赐、会盟和征伐缺一不可。

\term{后来改变了什么}：这种安排在王室强盛时扩大了周的影响；在王室衰弱时，封国积累的资源反而让地方成为独立力量。东周以后，齐、晋、楚、秦等大国的崛起都能追溯到这个更早的网络。

这项安排的次序比绝对年份更可靠。东征后的军事压力、对殷遗民的安置和成周的营建可以相互解释，却不能由此推定它们在某一年内依次完成。传世的\sourcebook{尚书·召诰}、\sourcebook{洛诰}保存了经营东方的王室语言，何尊留下较接近当时的王命记忆，\sourcebook{逸周书·作雒}则提供较完整而也更晚近的建城叙事。三类材料共同支持东方重组的轮廓，不能共同给出一张施工日期、迁置人数和封国边界都已精确无误的总图。

从现存材料看，周王室面对的选择不是“直接统治”或“完全放手”之间的二选一。它需要有一处较近的城市供会集和转运，也需要能在各地调动人手、粮食和祭祀关系的中介家族。这样，成周和封国便构成同一套安排的不同尺度：前者使王室能够抵达东方，后者使这种抵达不必每次都从关中重新开始；两者又都要靠地方社会持续承担劳役、供给和磨合的成本。
\end{event}

\begin{sourcenote}[成周到底留下了什么证据]
\sourcebook{尚书·召诰}、\sourcebook{尚书·洛诰}是营建成周的核心传世材料；何尊“余其宅兹中国”及其关于营建的追述，是最重要的出土文字之一。这里的“中国”指周初洛邑附近的政治中心，不是现代国家概念。洛阳地区的西周墓地与铸铜遗址，也让“东方安置”不只是后出的故事。王城与成周究竟是两座城、还是同一都市的不同功能区，仍有讨论。
\end{sourcenote}

\subsection{卫地：一段训诰怎样接住旧商核心}

\begin{event}{东征后的安置阶段（年月与路线不明）}{康叔受命经营卫地}{治理关联可信；具体迁置和地点不能坐实}{殷墟周边及卫地网络}\eventanchor{WZ-EARLY-WEI-FOUNDATION}{西周·卫地经营}
战事之后，最难的未必是再夺一座城，而是让命令留在旧商的核心地区。传统上与康叔相连的几篇训诰，把镜头落在这件并不显赫、却不能悬置的事上：谁来面对既有的人群和祭祀，谁来约束武力与刑罚，谁又能把王室的名分带进日常的治民事务。后世叙事把康叔受封、治卫和安置殷民串成一条清楚的故事线；从现存材料看，较能把握的是东征之后确有一项持续的地方经营，而不是一张注明日期、路线和户数的交接清单。

传世\sourcebook{尚书·康诰}所呈现的，是以德、刑与祭祀秩序相互牵连的治理语言。它没有留下一个现代意义的地方政府编制，却让人看见王室试图把惩戒的权限、对旧有习俗的警惕和受命者的责任放在同一套话语中。于是，\eventanchor{WZ-1040-KANG-SHU-RELOCATION}{西周·康叔治卫}所指的安置，不宜只理解为人口从甲地移到乙地：它还包括让新的受命家族、旧有人群、城邑劳役和祭祀权威能够在同一处政治关系里继续运作。

这里的资源不是抽象的“地方”。旧商核心区的人口、耕作与手工业基础，以及能够维持城邑和祭祀的物资，都是新秩序必须接住的条件；王室若没有可依托的地方中介，也难以持续监督这个区域。直接后果，是卫地被纳入周的宗法名分和军事—礼仪网络。中期的代价同样清楚：中心把治理交给有自己宗族根基的受命者，便须不断以册命、赏赐和裁断维系其联系，不能把一次受命误写成永久而无摩擦的服从。
\end{event}

\sourcebook{酒诰}把饮酒及其所象征的旧俗写成秩序风险，\eventanchor{WZ-EARLY-JIU-GAO}{西周·酒诰的训诫}；\sourcebook{梓材}又把筑城、用材和治民的话语并置，\eventanchor{WZ-EARLY-ZI-CAI}{西周·梓材的治民语言}。二者帮助理解周人如何命名治理难题，却不等于现场的禁酒令或工程日报。篇章针对的对象、颁布时点和今天所见文字是否保留一次完整训令，都不能确证。把它们同\eventanchor{WZ-EARLY-KANG-GAO}{西周·康诰的治理语言}放在一起，只能谨慎地说：早周统治者设想中的整合，既要有可执行的惩戒，也要有城邑、人力和祭祀能够承受的秩序。

三篇的差异也不应被一句“周初法令”抹平。《康诰》把刑罚、祭祀和受命者责任并置，《酒诰》把被归入旧俗的行为写作秩序风险，《梓材》则将筑城、用材与治民联系起来。它们提供的是相互呼应的治理语言，而非可逐条复原的地方条例。清华简《书》类材料带来的比较和传世篇章的层累问题，正要求读者区分：我们能够借它们看到王室如何设想整合，却不能确证每一篇都是一次原封不动保存下来的现场训令。

\begin{sourcenote}[谁在说卫地，谁没有留下声音]
接近事件的材料并不整齐。传世\sourcebook{尚书}诸篇保存了早周王命与训诰的政治语言，但其成篇、传抄和层累须逐篇辨析；殷墟的西周遗存能够显示旧都地区并未在征服后成为空白，却不能给出康叔行程或迁置人数。\sourcebook{史记·卫康叔世家}把康叔的经历编入完整世系，是汉代的后出叙事，适合与前两类材料互证，不宜当作现场记录。清华简《书》类材料与传世文本的比较，正促使研究者重新追问这些训诰的文本层次。现有证据尤其难以让我们听见被重新编组的殷民如何评价这套安排。
\end{sourcenote}

\subsection{燕地：城址、墓地与一支向北的网络}

\begin{event}{周初至成王时期说}{燕地封国中心逐步站稳}{始封传统较稳；始封者、绝对年代与墓主身份有争议}{燕山南麓的琉璃河一带}\eventanchor{WZ-EARLY-YAN-FOUNDATION}{西周·燕地封国中心}
向北的通道不能只靠谱系上的一次授封维持。一处地方中心若要站住脚，必须有可居住和守卫的城邑，有安葬并显示家族连续性的墓地，也有把受命关系交给后代的礼器。琉璃河的城址、墓地和燕侯器使这种缓慢的固定过程露出轮廓：这里不是一枚在地图上被涂上的颜色，而是地方精英、劳力、物资和北方交通被逐渐编接起来的节点。

这也解释了燕地在周初网络中的位置。王室需要能沟通北方通道、承担守卫与礼仪联系的中介；受命家族则需要把这种资格沉入城址、墓葬和器物所见的连续经营。其后果不是王室从此可以省去一切投入，而是多出一个能够汇集地方资源、也能够在代际之间保存自身声望的中心。这样的安排扩大了周的可达范围，同时也使地方能力有了积累的容器。
\end{event}

\begin{event}{西周早期（连续年代不明）}{琉璃河的城址、墓地与燕侯器}{考古遗存可靠；器主与墓主的对应关系仍有讨论}{燕地封国中心}\eventanchor{WZ-1030S-YAN-LIULIHE}{西周·琉璃河的燕侯器}
后来的\sourcebook{史记·燕召公世家}把燕的开端安放在召公家族的谱系里；真正落到地面的，却是城址的层位、墓地的排列和器物上的文字。燕侯旨鼎、燕侯盂与琉璃河的遗存并不能讲述某一天的授封仪式，却比一条整齐的世系更能说明：一个与周王室关系相连的精英中心，确实经历了持续的建设和安葬。

不过，器铭是为纪念、祭祀和家族传承而作，不是户籍册；墓葬的丰厚也不自动说明每一位墓主的政治身份。考古发掘报告提供的是同时代或接近同时代的物质证据，\sourcebook{史记}提供的是较晚而连贯的叙述，两者都不能单独裁定“第一位受封者是谁”或封国疆界到哪里。较稳妥的结论是，燕地中心的形成具有过程性；这使周的北方网络获得一处可依托的节点，也保留了具体人群、界线和年份尚未被材料填满的空白。
\end{event}

城址的建设、墓地的安排和礼器的制作，也使“封国”恢复了物质重量：它们需要劳力、原料和地方协作，并非王室册命自动产生。考古能够让这种投入的轮廓浮现，却难以分别识别每一位工匠、役夫和本地居民。燕侯器与琉璃河遗存因而既支持北方中心逐步稳定的判断，也要求叙述承认，最承担建设成本的人，往往最少以自己的语言出现在现有材料中。
